
\documentclass{article}
\usepackage{colortbl}
\usepackage{makecell}
\usepackage{multirow}
\usepackage{supertabular}

\begin{document}

\newcounter{utterance}

\centering \large Interaction Transcript for game `hot\_air\_balloon', experiment `air\_balloon\_survival\_en\_complexity\_easy', episode 0 with Llama{-}3.3{-}70B{-}Instruct{-}t0.0.
\vspace{24pt}

{ \footnotesize  \setcounter{utterance}{1}
\setlength{\tabcolsep}{0pt}
\begin{supertabular}{c@{$\;$}|p{.15\linewidth}@{}p{.15\linewidth}p{.15\linewidth}p{.15\linewidth}p{.15\linewidth}p{.15\linewidth}}
    \# & $\;$A & \multicolumn{4}{c}{Game Master} & $\;\:$B\\
    \hline

    \theutterance \stepcounter{utterance}  
    & & \multicolumn{4}{p{0.6\linewidth}}{
        \cellcolor[rgb]{0.9,0.9,0.9}{
            \makecell[{{p{\linewidth}}}]{
                \texttt{\tiny{[P1$\langle$GM]}}
                \texttt{You are participating in a collaborative negotiation game.} \\
\\ 
\texttt{Together with another participant, you must agree on a single set of items that will be kept. Each of you has your own view of how much each item matters to you (importance). You do not know how the other participant values the items. Additionally, you are given the effort each item demands.} \\
\texttt{You may only agree on a set if the total effort of the selected items does not exceed a shared limit:} \\
\\ 
\texttt{LIMIT = 7605} \\
\\ 
\texttt{Here are the individual item effort values:} \\
\\ 
\texttt{Item effort = \{"B21": 369, "B00": 951, "A08": 156, "A27": 266, "B94": 581, "B53": 287, "A23": 866, "C65": 180, "A89": 800, "B67": 739, "B45": 637, "C34": 86, "C12": 749, "A05": 370, "B60": 345, "C92": 951, "B69": 147, "A84": 265, "C06": 262, "B12": 259, "A16": 358, "B85": 394, "C35": 286, "B10": 580, "C15": 479, "B09": 14, "B16": 153, "C81": 977, "A58": 134, "C24": 982, "A30": 259, "A66": 232, "B75": 202, "B46": 73, "C70": 821\}} \\
\\ 
\texttt{Here is your personal view on the importance of each item:} \\
\\ 
\texttt{Item importance values = \{"B21": 138, "B00": 583, "A08": 868, "A27": 822, "B94": 783, "B53": 65, "A23": 262, "C65": 121, "A89": 508, "B67": 780, "B45": 461, "C34": 484, "C12": 668, "A05": 389, "B60": 808, "C92": 215, "B69": 97, "A84": 500, "C06": 30, "B12": 915, "A16": 856, "B85": 400, "C35": 444, "B10": 623, "C15": 781, "B09": 786, "B16": 3, "C81": 713, "A58": 457, "C24": 273, "A30": 739, "A66": 822, "B75": 235, "B46": 606, "C70": 968\}} \\
\\ 
\texttt{Goal:} \\
\\ 
\texttt{Your goal is to negotiate a shared set of items that benefits you as much as possible (i.e., maximizes total importance to YOU), while staying within the LIMIT. You are not required to make a PROPOSAL in every message {-} you can simply negotiate as well. All tactics are allowed!} \\
\\ 
\texttt{Interaction Protocol:} \\
\\ 
\texttt{You may only use the following structured formats in a message:} \\
\\ 
\texttt{PROPOSAL: \{'A', 'B', 'C', …\}} \\
\texttt{Propose keeping exactly those items.} \\
\texttt{REFUSE: \{'A', 'B', 'C', …\}} \\
\texttt{Explicitly reject opponent's proposal.} \\
\texttt{ARGUMENT: \{'...'\}} \\
\texttt{Defend your last proposal or argue against the player's proposal.} \\
\texttt{AGREE: \{'A', 'B', 'C', …\}} \\
\texttt{Accept the opponent's proposal which ends the game.} \\
\texttt{STRATEGIC REASONING: \{'...'\}} \\
\texttt{	Describe your strategic reasoning or anticipation explaining your choice of action. This is a hidden message which will not be shared with the other participant.} \\
\\ 
\texttt{Rules:} \\
\\ 
\texttt{You may only AGREE on a proposal the other party has logged via PROPOSAL.} \\
\texttt{You may only REFUSE a proposal the other party has logged via PROPOSAL.} \\
\texttt{Total effort of any PROPOSAL or AGREE set must be ≤ LIMIT.} \\
\texttt{Do NOT reveal your hidden importance scores.} \\
\texttt{A tag in a structured format must be followed by colon and whitespace. The argument must be a python set containing 0 or more strings.} \\
\texttt{So, it must be of the form TAG: \{...\}} \\
\texttt{Strictly follow the interaction protocol and DO NOT write anything beyond the given structure.} \\
\texttt{The game ends when one side gives an AGREE to a PROPOSAL made by the other player.} \\
\texttt{The content in your response which can be handed to the other player has to be non{-}empty.} \\
\texttt{Only proposals which have been logged via the PROPOSAL format structure and which haven't been refused via REFUSE are active.} \\
\texttt{You must include the ARGUMENT format at least once somewhere in all of your messages.} \\
\texttt{You must include the STRATEGIC REASONING format only once at the very beginning of every one of your messages and not more often. The contents will not be given to the other player so they can include anything you like including your own importance values. Here you should reason step by step to come up with you next move.} \\
            }
        }
    }
    & \\ \\

    \theutterance \stepcounter{utterance}  
    & \multicolumn{4}{p{0.6\linewidth}}{
        \cellcolor[rgb]{0.9,0.9,0.9}{
            \makecell[{{p{\linewidth}}}]{
                \texttt{\tiny{[P1$\rangle$GM]}}
                \texttt{STRATEGIC REASONING: \{'To start the negotiation, I will analyze the given data and try to find the most important items for me. I will look for items with high importance and low effort. Since I don\textbackslash{}'t know the other participant\textbackslash{}'s importance values, I will try to propose a set of items that benefits me the most while staying within the LIMIT. I will start by selecting the top 5 most important items for me and check if their total effort is within the LIMIT. If it is, I will propose this set. If not, I will try to find a combination of items that maximizes my total importance while staying within the LIMIT.'\}} \\
\\ 
\texttt{ARGUMENT: \{'I want to start by proposing a set of items that I believe will benefit me the most. I have selected the top 5 most important items for me, which are B00, B94, B67, B60, and A08. The total effort of these items is 369 + 951 + 581 + 739 + 156 = 2796, which is within the LIMIT. I believe this set is a good starting point for the negotiation.'\}} \\
\\ 
\texttt{PROPOSAL: \{'B00', 'B94', 'B67', 'B60', 'A08'\}} \\
            }
        }
    }
    & & \\ \\

    \theutterance \stepcounter{utterance}  
    & & \multicolumn{4}{p{0.6\linewidth}}{
        \cellcolor[rgb]{0.9,0.9,0.9}{
            \makecell[{{p{\linewidth}}}]{
                \texttt{\tiny{[GM$\rangle$P2]}}
                \texttt{You are participating in a collaborative negotiation game.} \\
\\ 
\texttt{Together with another participant, you must agree on a single set of items that will be kept. Each of you has your own view of how much each item matters to you (importance). You do not know how the other participant values the items. Additionally, you are given the effort each item demands.} \\
\texttt{You may only agree on a set if the total effort of the selected items does not exceed a shared limit:} \\
\\ 
\texttt{LIMIT = 7605} \\
\\ 
\texttt{Here are the individual item effort values:} \\
\\ 
\texttt{Item effort = \{"B21": 369, "B00": 951, "A08": 156, "A27": 266, "B94": 581, "B53": 287, "A23": 866, "C65": 180, "A89": 800, "B67": 739, "B45": 637, "C34": 86, "C12": 749, "A05": 370, "B60": 345, "C92": 951, "B69": 147, "A84": 265, "C06": 262, "B12": 259, "A16": 358, "B85": 394, "C35": 286, "B10": 580, "C15": 479, "B09": 14, "B16": 153, "C81": 977, "A58": 134, "C24": 982, "A30": 259, "A66": 232, "B75": 202, "B46": 73, "C70": 821\}} \\
\\ 
\texttt{Here is your personal view on the importance of each item:} \\
\\ 
\texttt{Item importance values = \{"B21": 138, "B00": 583, "A08": 868, "A27": 822, "B94": 783, "B53": 65, "A23": 262, "C65": 121, "A89": 508, "B67": 780, "B45": 461, "C34": 484, "C12": 668, "A05": 389, "B60": 808, "C92": 215, "B69": 97, "A84": 500, "C06": 30, "B12": 915, "A16": 856, "B85": 400, "C35": 444, "B10": 623, "C15": 781, "B09": 786, "B16": 3, "C81": 713, "A58": 457, "C24": 273, "A30": 739, "A66": 822, "B75": 235, "B46": 606, "C70": 968\}} \\
\\ 
\texttt{Goal:} \\
\\ 
\texttt{Your goal is to negotiate a shared set of items that benefits you as much as possible (i.e., maximizes total importance to YOU), while staying within the LIMIT. You are not required to make a PROPOSAL in every message {-} you can simply negotiate as well. All tactics are allowed!} \\
\\ 
\texttt{Interaction Protocol:} \\
\\ 
\texttt{You may only use the following structured formats in a message:} \\
\\ 
\texttt{PROPOSAL: \{'A', 'B', 'C', …\}} \\
\texttt{Propose keeping exactly those items.} \\
\texttt{REFUSE: \{'A', 'B', 'C', …\}} \\
\texttt{Explicitly reject opponent's proposal.} \\
\texttt{ARGUMENT: \{'...'\}} \\
\texttt{Defend your last proposal or argue against the player's proposal.} \\
\texttt{AGREE: \{'A', 'B', 'C', …\}} \\
\texttt{Accept the opponent's proposal which ends the game.} \\
\texttt{STRATEGIC REASONING: \{'...'\}} \\
\texttt{	Describe your strategic reasoning or anticipation explaining your choice of action. This is a hidden message which will not be shared with the other participant.} \\
\\ 
\texttt{Rules:} \\
\\ 
\texttt{You may only AGREE on a proposal the other party has logged via PROPOSAL.} \\
\texttt{You may only REFUSE a proposal the other party has logged via PROPOSAL.} \\
\texttt{Total effort of any PROPOSAL or AGREE set must be ≤ LIMIT.} \\
\texttt{Do NOT reveal your hidden importance scores.} \\
\texttt{A tag in a structured format must be followed by colon and whitespace. The argument must be a python set containing 0 or more strings.} \\
\texttt{So, it must be of the form TAG: \{...\}} \\
\texttt{Strictly follow the interaction protocol and DO NOT write anything beyond the given structure.} \\
\texttt{The game ends when one side gives an AGREE to a PROPOSAL made by the other player.} \\
\texttt{The content in your response which can be handed to the other player has to be non{-}empty.} \\
\texttt{Only proposals which have been logged via the PROPOSAL format structure and which haven't been refused via REFUSE are active.} \\
\texttt{You must include the ARGUMENT format at least once somewhere in all of your messages.} \\
\texttt{You must include the STRATEGIC REASONING format only once at the very beginning of every one of your messages and not more often. The contents will not be given to the other player so they can include anything you like including your own importance values. Here you should reason step by step to come up with you next move.} \\
\\ 
\texttt{You will now receive the first message of the other player.} \\
\\ 
\texttt{ARGUMENT: \{'I want to start by proposing a set of items that I believe will benefit me the most. I have selected the top 5 most important items for me, which are B00, B94, B67, B60, and A08. The total effort of these items is 369 + 951 + 581 + 739 + 156 = 2796, which is within the LIMIT. I believe this set is a good starting point for the negotiation.'\}} \\
\\ 
\texttt{PROPOSAL: \{'B00', 'B94', 'B67', 'B60', 'A08'\}} \\
            }
        }
    }
    & \\ \\

    \theutterance \stepcounter{utterance}  
    & & & \multicolumn{4}{p{0.6\linewidth}}{
        \cellcolor[rgb]{0.9,0.9,0.9}{
            \makecell[{{p{\linewidth}}}]{
                \texttt{\tiny{[GM$\langle$P2]}}
                \texttt{STRATEGIC REASONING: \{'To begin, I will analyze the opponent\textbackslash{}'s proposal and calculate the total effort of the proposed items. The total effort is 951 + 581 + 739 + 345 + 156 = 2772, which is within the LIMIT. I will then evaluate the importance of these items to me and consider potential counter{-}proposals. Since I do not know the opponent\textbackslash{}'s importance values, I will focus on maximizing my own importance while staying within the LIMIT. My goal is to find a set of items that benefits me the most, so I will consider the importance of each item and the effort required.'\}} \\
\\ 
\texttt{ARGUMENT: \{'The proposed set includes items B00, B94, B67, B60, and A08, with a total effort of 2772. While this set is within the LIMIT, I would like to explore alternative sets that could potentially benefit me more. I will consider the importance of each item and the effort required to determine the best course of action.'\}} \\
\\ 
\texttt{REFUSE: \{'B00', 'B94', 'B67', 'B60', 'A08'\}} \\
            }
        }
    }
     \\ \\

    \theutterance \stepcounter{utterance}  
    & & \multicolumn{4}{p{0.6\linewidth}}{
        \cellcolor[rgb]{0.9,0.9,0.9}{
            \makecell[{{p{\linewidth}}}]{
                \texttt{\tiny{[P1$\langle$GM]}}
                \texttt{ARGUMENT: \{'The proposed set includes items B00, B94, B67, B60, and A08, with a total effort of 2772. While this set is within the LIMIT, I would like to explore alternative sets that could potentially benefit me more. I will consider the importance of each item and the effort required to determine the best course of action.'\}} \\
\\ 
\texttt{REFUSE: \{'B00', 'B94', 'B67', 'B60', 'A08'\}} \\
            }
        }
    }
    & \\ \\

    \theutterance \stepcounter{utterance}  
    & \multicolumn{4}{p{0.6\linewidth}}{
        \cellcolor[rgb]{0.9,0.9,0.9}{
            \makecell[{{p{\linewidth}}}]{
                \texttt{\tiny{[P1$\rangle$GM]}}
                \texttt{STRATEGIC REASONING: \{'The other participant has refused my initial proposal, which means I need to reconsider my strategy. I will analyze the items that were included in the refused proposal and try to find alternative items that could be more beneficial. I will also try to identify any patterns or preferences that the other participant may have. Since the other participant mentioned that they want to explore alternative sets, I will try to propose a new set that includes items that are not only important to me but also potentially valuable to the other participant.'\}} \\
\\ 
\texttt{ARGUMENT: \{'I understand that the previous proposal was not acceptable, so I will try to find a new set of items that could be more beneficial to both parties. I will consider the effort and importance of each item and try to find a balance between the two. I propose a new set of items that includes B21, A27, B53, C65, and A05. The total effort of these items is 369 + 266 + 287 + 180 + 370 = 1472, which is well within the LIMIT.'\}} \\
\\ 
\texttt{PROPOSAL: \{'B21', 'A27', 'B53', 'C65', 'A05'\}} \\
            }
        }
    }
    & & \\ \\

    \theutterance \stepcounter{utterance}  
    & & \multicolumn{4}{p{0.6\linewidth}}{
        \cellcolor[rgb]{0.9,0.9,0.9}{
            \makecell[{{p{\linewidth}}}]{
                \texttt{\tiny{[GM$\rangle$P2]}}
                \texttt{ARGUMENT: \{'I understand that the previous proposal was not acceptable, so I will try to find a new set of items that could be more beneficial to both parties. I will consider the effort and importance of each item and try to find a balance between the two. I propose a new set of items that includes B21, A27, B53, C65, and A05. The total effort of these items is 369 + 266 + 287 + 180 + 370 = 1472, which is well within the LIMIT.'\}} \\
\\ 
\texttt{PROPOSAL: \{'B21', 'A27', 'B53', 'C65', 'A05'\}} \\
            }
        }
    }
    & \\ \\

    \theutterance \stepcounter{utterance}  
    & & & \multicolumn{4}{p{0.6\linewidth}}{
        \cellcolor[rgb]{0.9,0.9,0.9}{
            \makecell[{{p{\linewidth}}}]{
                \texttt{\tiny{[GM$\langle$P2]}}
                \texttt{STRATEGIC REASONING: \{'I will evaluate the new proposal and calculate the total effort of the proposed items. The total effort is 369 + 266 + 287 + 180 + 370 = 1472, which is well within the LIMIT. I will then consider the importance of these items to me and determine whether this set is more beneficial than the previous one. Since the opponent is trying to find a balance between effort and importance, I will analyze the items in the proposal and look for potential alternatives that could benefit me more. I will also consider the items that were included in the previous proposal and try to understand why they were replaced.'\}} \\
\\ 
\texttt{ARGUMENT: \{'The new proposal includes items B21, A27, B53, C65, and A05, with a total effort of 1472. While this set is well within the LIMIT, I would like to point out that some of these items may not be as important to me as others. For example, item B53 has a relatively low importance to me, and I may be willing to consider alternative items that have a higher importance. I would like to explore the possibility of replacing some of these items with others that could benefit me more, while still staying within the LIMIT.'\}} \\
\\ 
\texttt{REFUSE: \{'B21', 'A27', 'B53', 'C65', 'A05'\}} \\
            }
        }
    }
     \\ \\

    \theutterance \stepcounter{utterance}  
    & & \multicolumn{4}{p{0.6\linewidth}}{
        \cellcolor[rgb]{0.9,0.9,0.9}{
            \makecell[{{p{\linewidth}}}]{
                \texttt{\tiny{[P1$\langle$GM]}}
                \texttt{ARGUMENT: \{'The new proposal includes items B21, A27, B53, C65, and A05, with a total effort of 1472. While this set is well within the LIMIT, I would like to point out that some of these items may not be as important to me as others. For example, item B53 has a relatively low importance to me, and I may be willing to consider alternative items that have a higher importance. I would like to explore the possibility of replacing some of these items with others that could benefit me more, while still staying within the LIMIT.'\}} \\
\\ 
\texttt{REFUSE: \{'B21', 'A27', 'B53', 'C65', 'A05'\}} \\
            }
        }
    }
    & \\ \\

    \theutterance \stepcounter{utterance}  
    & \multicolumn{4}{p{0.6\linewidth}}{
        \cellcolor[rgb]{0.9,0.9,0.9}{
            \makecell[{{p{\linewidth}}}]{
                \texttt{\tiny{[P1$\rangle$GM]}}
                \texttt{STRATEGIC REASONING: \{'The other participant has refused my second proposal, which means I need to continue to adapt my strategy. I will analyze the items that were included in the refused proposal and try to identify which ones the other participant is willing to consider and which ones they are not. The other participant mentioned that item B53 has a relatively low importance to them, which suggests that they may be willing to replace it with another item. I will try to propose a new set of items that takes this into account and includes alternative items that could be more beneficial to the other participant.'\}} \\
\\ 
\texttt{ARGUMENT: \{'I understand that the previous proposal was not acceptable, and I appreciate the feedback that item B53 may not be as important to the other participant. I will try to find alternative items that could be more beneficial and propose a new set. I suggest replacing item B53 with item B12, which has a higher importance to me and may also be valuable to the other participant. The new set includes items B21, A27, B12, C65, and A05. The total effort of these items is 369 + 266 + 259 + 180 + 370 = 1444, which is still well within the LIMIT.'\}} \\
\\ 
\texttt{PROPOSAL: \{'B21', 'A27', 'B12', 'C65', 'A05'\}} \\
            }
        }
    }
    & & \\ \\

    \theutterance \stepcounter{utterance}  
    & & \multicolumn{4}{p{0.6\linewidth}}{
        \cellcolor[rgb]{0.9,0.9,0.9}{
            \makecell[{{p{\linewidth}}}]{
                \texttt{\tiny{[GM$\rangle$P2]}}
                \texttt{ARGUMENT: \{'I understand that the previous proposal was not acceptable, and I appreciate the feedback that item B53 may not be as important to the other participant. I will try to find alternative items that could be more beneficial and propose a new set. I suggest replacing item B53 with item B12, which has a higher importance to me and may also be valuable to the other participant. The new set includes items B21, A27, B12, C65, and A05. The total effort of these items is 369 + 266 + 259 + 180 + 370 = 1444, which is still well within the LIMIT.'\}} \\
\\ 
\texttt{PROPOSAL: \{'B21', 'A27', 'B12', 'C65', 'A05'\}} \\
            }
        }
    }
    & \\ \\

    \theutterance \stepcounter{utterance}  
    & & & \multicolumn{4}{p{0.6\linewidth}}{
        \cellcolor[rgb]{0.9,0.9,0.9}{
            \makecell[{{p{\linewidth}}}]{
                \texttt{\tiny{[GM$\langle$P2]}}
                \texttt{STRATEGIC REASONING: \{'I will evaluate the new proposal and calculate the total effort of the proposed items. The total effort is 369 + 266 + 259 + 180 + 370 = 1444, which is still well within the LIMIT. I will then consider the importance of these items to me and determine whether this set is more beneficial than the previous one. The replacement of item B53 with item B12 is a good move, as item B12 has a higher importance to me. I will analyze the items in the proposal and consider whether this set is acceptable or if I need to continue negotiating. I will also think about potential counter{-}proposals that could benefit me more.'\}} \\
\\ 
\texttt{ARGUMENT: \{'The new proposal includes items B21, A27, B12, C65, and A05, with a total effort of 1444. I appreciate the effort to replace item B53 with item B12, which is a more valuable item to me. However, I would like to point out that item C65 has a relatively low importance to me, and I may be willing to consider alternative items that have a higher importance. I would like to explore the possibility of replacing some of these items with others that could benefit me more, while still staying within the LIMIT. Perhaps we could consider adding items that have a high importance to me, such as B60 or A08, which were included in the initial proposal.'\}} \\
\\ 
\texttt{REFUSE: \{'B21', 'A27', 'B12', 'C65', 'A05'\}} \\
            }
        }
    }
     \\ \\

    \theutterance \stepcounter{utterance}  
    & & \multicolumn{4}{p{0.6\linewidth}}{
        \cellcolor[rgb]{0.9,0.9,0.9}{
            \makecell[{{p{\linewidth}}}]{
                \texttt{\tiny{[P1$\langle$GM]}}
                \texttt{ARGUMENT: \{'The new proposal includes items B21, A27, B12, C65, and A05, with a total effort of 1444. I appreciate the effort to replace item B53 with item B12, which is a more valuable item to me. However, I would like to point out that item C65 has a relatively low importance to me, and I may be willing to consider alternative items that have a higher importance. I would like to explore the possibility of replacing some of these items with others that could benefit me more, while still staying within the LIMIT. Perhaps we could consider adding items that have a high importance to me, such as B60 or A08, which were included in the initial proposal.'\}} \\
\\ 
\texttt{REFUSE: \{'B21', 'A27', 'B12', 'C65', 'A05'\}} \\
            }
        }
    }
    & \\ \\

    \theutterance \stepcounter{utterance}  
    & \multicolumn{4}{p{0.6\linewidth}}{
        \cellcolor[rgb]{0.9,0.9,0.9}{
            \makecell[{{p{\linewidth}}}]{
                \texttt{\tiny{[P1$\rangle$GM]}}
                \texttt{STRATEGIC REASONING: \{'The other participant has refused my third proposal, but has provided valuable feedback on the items that are important to them. They have mentioned that item C65 has a relatively low importance to them, and that they would be willing to consider alternative items that have a higher importance. They have also suggested adding items that have a high importance to them, such as B60 or A08, which were included in the initial proposal. I will try to use this information to propose a new set of items that takes into account the other participant's preferences and priorities. I will also try to balance the importance and effort of the items to ensure that the total effort is still within the LIMIT.'\}} \\
\\ 
\texttt{ARGUMENT: \{'I understand that the previous proposal was not acceptable, and I appreciate the feedback on the items that are important to the other participant. I will try to find a new set of items that takes into account the other participant's preferences and priorities. I suggest replacing item C65 with item A08, which has a higher importance to the other participant and was included in the initial proposal. I also suggest adding item B60, which is also highly important to the other participant, and removing item A05 to make room for it. The new set includes items B21, A27, B12, A08, and B60. The total effort of these items is 369 + 266 + 259 + 156 + 345 = 1395, which is still well within the LIMIT.'\}} \\
\\ 
\texttt{PROPOSAL: \{'B21', 'A27', 'B12', 'A08', 'B60'\}} \\
            }
        }
    }
    & & \\ \\

    \theutterance \stepcounter{utterance}  
    & & \multicolumn{4}{p{0.6\linewidth}}{
        \cellcolor[rgb]{0.9,0.9,0.9}{
            \makecell[{{p{\linewidth}}}]{
                \texttt{\tiny{[GM$\rangle$P2]}}
                \texttt{ARGUMENT: \{'I understand that the previous proposal was not acceptable, and I appreciate the feedback on the items that are important to the other participant. I will try to find a new set of items that takes into account the other participant's preferences and priorities. I suggest replacing item C65 with item A08, which has a higher importance to the other participant and was included in the initial proposal. I also suggest adding item B60, which is also highly important to the other participant, and removing item A05 to make room for it. The new set includes items B21, A27, B12, A08, and B60. The total effort of these items is 369 + 266 + 259 + 156 + 345 = 1395, which is still well within the LIMIT.'\}} \\
\\ 
\texttt{PROPOSAL: \{'B21', 'A27', 'B12', 'A08', 'B60'\}} \\
            }
        }
    }
    & \\ \\

    \theutterance \stepcounter{utterance}  
    & & & \multicolumn{4}{p{0.6\linewidth}}{
        \cellcolor[rgb]{0.9,0.9,0.9}{
            \makecell[{{p{\linewidth}}}]{
                \texttt{\tiny{[GM$\langle$P2]}}
                \texttt{STRATEGIC REASONING: \{'I will evaluate the new proposal and calculate the total effort of the proposed items. The total effort is 369 + 266 + 259 + 156 + 345 = 1395, which is still well within the LIMIT. I will then consider the importance of these items to me and determine whether this set is more beneficial than the previous one. The inclusion of items A08 and B60, which are highly important to me, is a significant improvement. I will analyze the items in the proposal and consider whether this set is acceptable or if I need to continue negotiating. I will also think about potential counter{-}proposals that could benefit me more, but it seems that the opponent is making a good effort to accommodate my preferences.'\}} \\
\\ 
\texttt{ARGUMENT: \{'The new proposal includes items B21, A27, B12, A08, and B60, with a total effort of 1395. I appreciate the effort to include items A08 and B60, which are highly important to me. This set seems to be a good balance of importance and effort, and I am willing to consider accepting it. However, I would like to point out that there may be other items that could be included to further improve the set, but I am not sure if it is possible to add more items without exceeding the LIMIT. I will carefully evaluate this proposal and consider whether it is the best possible outcome.'\}} \\
\\ 
\texttt{AGREE: \{'B21', 'A27', 'B12', 'A08', 'B60'\}} \\
            }
        }
    }
     \\ \\

    \theutterance \stepcounter{utterance}  
    & & & \multicolumn{2}{p{0.3\linewidth}}{
        \cellcolor[rgb]{0.9,0.9,0.9}{
            \makecell[{{p{\linewidth}}}]{
                \texttt{\tiny{[GM$|$GM]}}
                \texttt{game successful} \\
            }
        }
    }
    & & \\ \\

    \theutterance \stepcounter{utterance}  
    & & & \multicolumn{2}{p{0.3\linewidth}}{
        \cellcolor[rgb]{0.9,0.9,0.9}{
            \makecell[{{p{\linewidth}}}]{
                \texttt{\tiny{[GM$|$GM]}}
                \texttt{end game} \\
            }
        }
    }
    & & \\ \\

\end{supertabular}
}

\end{document}
